%% kapa: exam=1 round=36 year=2025 period=1 form=A subject=economics subject_ko=경제학원론 seq=37 number=77
%% source: Q-Net 감정평가사 기출문제 (cst00309) artlSeq=5231525 file=raw/1008/5231525/2025년 제36회 감정평가사 1차시험 1교시 문제지 원본.pdf
%% conversion: kordoc 4.24.51 -> tools/gen1.py (auto); review=proofread-sonnet; flags=none
\begin{question}{77}
아래의 거시경제모형의 균형에서 균형임금보다 5를 증가시킨 임금 수준의 최저임금제를 도입할 경우, 거시경제 변수의 변화에 관한 설명으로 옳은 것은? (단, $Y$는 총생산, $r$은 실질이자율, $W$는 명목임금)
\begin{framedtext}
$\circ$ 총저축: $S=s(r)Y$ \par
$\circ$ 저축율: $s(r)=0.01r$ \par
$\circ$ 투자: $I=300-r$ \par
$\circ$ 총생산: $Y=2L$ \par
$\circ$ 노동수요: $L^{d}=300-10\frac{W}{P}$ \par
$\circ$ 노동공급: $L^{s}=10\frac{W}{P}$ \par
$\circ$ 노동량: $L=L^{s}=L^{d}$ (균형일 경우) \par
$\circ$ 노동량: $L=L^{d}$ (불균형일 경우) \par
$\circ$ 소비: $C=Y-S$ \par
$\circ$ 물가: $P=1$
\end{framedtext}
\begin{choices}
  \item 총생산은 증가한다.
  \item 이자율은 하락한다.
  \item 총저축은 증가한다.
  \item 고용량은 증가한다.
  \item 소비는 감소한다.
\end{choices}
\end{question}
\answer{5}
